%%% Choose between 16:9 and 4:3 format by commenting out/uncommenting one of the following lines:
\documentclass[aspectratio=169]{beamer} % 16:9
% \documentclass{beamer} % 4:3

%=========================================================================================================================

\usepackage{beamerthemesplit}
\usepackage{appendixnumberbeamer}
\usepackage[english]{babel}
\usepackage[latin1]{inputenc}

\usepackage{tikz}
\usepackage{url}
\usepackage{tabularx}

\useinnertheme{aig}
\useoutertheme{aig}
\usecolortheme{aig}
\usefonttheme{aig}


%=========================================================================================================================
\title{Introduction into Possibilistic Logic}
\subtitle{Possibilistic Logic - An Overview by Didier Dubois and Henri Prade}
\author[A. K�hn]{Alexander K�hn}
\institute{Artificial Intelligence Group,\\
University of Hagen, Germany}
\date{December 17, 2024}
%=========================================================================================================================
\logo{\includegraphics[width=3cm]{figures/logoaig.png}}
%=========================================================================================================================

\begin{document}

%=========================================================================================================================

\begin{frame}
  \titlepage
\end{frame}
\nologo

% Greeting
% Welcome to my presentation about Possibilistic Logic
% I am currently in my second year in the Master Praktische Informatik
% As a main source I have used the overview chapter: Possibilistic Logic - An Overview Didier Dubois and Henri Prade

% Overview
% After a short Introduction and history,
% we will dive into possibility theory, as this is necessary to step further to possibilitc logic
% I will introduce you to possibilistic logic and explain the semantic and syntactic aspects
% Following with the explanation, I am providing an example
% and end with information about real world usage


\section{Introduction}
% Slide 1: Introduction
\begin{frame}{What is Possibilistic Logic?}
	\begin{itemize}
		\item \textbf{Definition}: Possibilistic logic is a weighted logic that handles uncertainty in a qualitative way by adding certainty to classical logic formulas		
		\item \textbf{Applications}: Used in expert systems, decision-making, and AI for reasoning under uncertainty.
	\end{itemize}
\end{frame}
% ability to deal with inconsistency
% modal flavor
% remains very close to classical logic
% developed for about thirty years
% it can also model preferences

\section{Possibility Theory}
% Now I will introduce you first into Possibilty theory
\begin{frame}{What are Possibility Distributions?}
	\begin{itemize}
		\item A \textbf{possibility distribution} is a mapping:
		\[
		\pi: U \to S,
		\]
		where:
		\begin{itemize}
			\item \( U \): Set of states
			\item \( S \): Totally ordered scale (\( S = [0, 1] \) or finite chains).
			\item \( \pi(u) \): Degree of possibility (plausibility) of state \( u \).
		\end{itemize}
		\item \( \pi(u) = 0 \): State \( u \) is \textbf{impossible}.
		\item \( \pi(u) = 1 \): State \( u \) is \textbf{fully possible}.
	\end{itemize}
	\begin{block}{Normalization Condition}
		At least one state \( u \in U \) must have \( \pi(u) = 1 \) (an actual state exists).
	\end{block}
\end{frame}
% To start with possibility theory we need to first understand possibility distributions
% A possibility distribution is a mapping ? from U to a totally ordered scale S, with top denoted by 1 and bottom by 0
%function ? represents the state of knowledge of an agent (about the actual state of a?airs), also called an epistemic state distinguishing what is plausible from what is less plausible, what is the normal course of things from what is not, what is surprising from what is expected. It represents a flexible restriction on what is the actual state of facts with the following conventions (similar to probability

\begin{frame}{Specificity of a Possibility Distribution}
	\begin{itemize}
		\item A distribution \( \pi \) is \textbf{at least as specific} as \( \pi' \) if:
		\[
		\pi(u) \leq \pi'(u), \quad \forall u \in U.
		\]
		\item Examples:
		\begin{itemize}
			\item \textbf{Complete knowledge:} \( \pi(u_0) = 1 \), \( \pi(u) = 0 \) for all \( u \neq u_0 \).
			\item \textbf{Complete ignorance:} \( \pi(u) = 1 \) for all \( u \in U \).
		\end{itemize}
		\item Principle of minimal specificity: It states that any hypothesis not known to be impossible cannot be ruled out. It is a minimal commitment, cautious, information principle.
	\end{itemize}
\end{frame}
% A possibility distribution ? is said to be at least as specific as another ??if and only if for each state of a?airs u: ?(u) ???(u) Then, ? is at least as restrictive and informative as ??, since it rules out at least as many states with at least as much strength.

\begin{frame}{Possibility and Necessity Functions}
	\begin{itemize}
		\item Possibility measure:
		\[
		\Pi(A) = \sup_{u \in A} \pi(u).
		\]
		Evaluates the extent to which \( A \) is consistent with \( \pi \).
		\item Necessity measure:
		\[
		N(A) = \inf_{u \notin A} (1 - \pi(u)).
		\]
		Evaluates the extent to which \( A \) is certainly implied by \( \pi \).
	\end{itemize}
	\begin{block}{Duality Relationship}
		\[
		N(A) = 1 - \Pi(A^c).
		\]
		Where \( A^c \) is the complement of \( A \).
	\end{block}
\end{frame}
% Given a simple query of the form ?does event A occur?? (is the corresponding proposition a true?), where A is a subset of states, the set of models of a, the response to the query can be obtained by computing degrees of possibility and necessity respectively


\begin{frame}{Key Properties of Measures}
	\begin{itemize}
		\item \textbf{Possibility: Maxitivity Property}
		\[
		\Pi(A \cup B) = \max(\Pi(A), \Pi(B)).
		\]
		\item \textbf{Necessity: Minimization Property}
		\[
		N(A \cap B) = \min(N(A), N(B)).
		\]
	\end{itemize}
\end{frame}
% Possibility measures satisfy the characteristic ?maxitivity? property
% Necessity measures satisfy an axiom dual to that of possibility measures
% N(a?b) = min(N(a),N(b)), expresses that to be certain about a?b, one has to be certain about a and to be certain about b.
% Note also that we only have N(A?B) ? max(N(A),N(B)).This goes well with the idea that one may be certain about the event A?B, without being really certain about more specific events such as A and B.

\begin{frame}{Certainty Qualification in Possibilistic Logic}
	\begin{itemize}
		\item Certainty-qualified information: "A is certain to degree \(\alpha\)" is modeled by the constraint: 
		\[
		N(A) \geq \alpha
		\]
		\item Least Specific Possibility Distribution:
		\[
		\pi_{(A, \alpha)}(u) = 
		\begin{cases} 
			1 & \text{if } u \in A, \\ 
			1 - \alpha & \text{otherwise}.
		\end{cases}
		\]
		\begin{itemize}
			\item If \(\alpha = 1\): \(\pi(A, \alpha)\) becomes the characteristic function of \(A\).
			\item If \(\alpha = 0\): \(\pi(A, \alpha)\) represents total ignorance.
		\end{itemize}
	\end{itemize}
\end{frame}
% Human knowledge is often expressed as statements with **belief degrees**.
%  Represents a family of possible epistemic states \(\pi\) that satisfy this constraint.
% Serves as a building block for constructing possibility distributions from multiple uncertain knowledge pieces.
% Integral to the semantics of possibilistic logic.

\begin{frame}{Example}
	\begin{itemize}
		\item A doctor is trying to diagnose a patient based on symptoms and diseases:
		\begin{itemize}
			\item \( U = \{\text{flu}, \text{cold}, \text{allergy}\} \).
			\item Possibility distribution \(\pi\) based on symptom severity:
			\[
			\pi(\text{flu}) = 0.9, \quad \pi(\text{cold}) = 0.6, \quad \pi(\text{allergy}) = 0.3.
			\]
		\end{itemize}
		\item The doctor evaluates the possibility and necessity measures for a hypothesis \( A \):
		\[ 
		A = \{\text{flu}, \text{cold}\}
		\]
		\item Calculation:
		\begin{align*}
			\Pi(A) &= \sup_{u \in A} \pi(u) = \max(\pi(\text{flu}), \pi(\text{cold})) = 0.9. \\
			N(A) &= \inf_{u \notin A} (1 - \pi(u)) = 1 - \pi(\text{allergy}) = 1 - 0.3 = 0.7.
		\end{align*}
		\item Interpretation:
		\begin{itemize}
			\item Possibility (\(\Pi(A)\)): Flu or cold is highly possible (\(0.9\)).
			\item Necessity (\(N(A)\)): It is moderately certain (\(0.7\)) that the patient has either flu or cold.
		\end{itemize}
	\end{itemize}
\end{frame}



% Slide 2: Syntax
\section{Syntax of Possibilistic Logic}
\begin{frame}{Possibilistic Logic}
	\begin{itemize}
		\item A basic  possibilistic logic formula is represented as:
		\[
		(a, \alpha)
		\]
		where:
		\begin{itemize}
			\item \(a\): A classical logic formula (e.g., \(p \to q\)).
			\item \(\alpha \in [0, 1]\): Certainty level (minimum necessity degree of \(a\)).
		\end{itemize}
		\item Interpretation:
		\begin{itemize}
			\item \( (a, \alpha) \): "Formula \(a\) is at least certain to degree \(\alpha\)."
			\item The higher \(\alpha\), the stronger the belief in \(a\).
		\end{itemize}
	\end{itemize}
\end{frame}
% A basic possibilistic logic formula is a pair (a,?) made of a classical logic formula a associated with a certainty level ? ?(0,1], viewed as a lower bound of a necessity measure, i.e., (a,?) is understood as N(a) ??
% In the following, we only consider the case of (basic) possibilistic propositional logic, ?L for short, i.e., possibilistic logic formulas (a,?) are such that a is a formula in a propositional language

% Slide 4: Inference
\begin{frame}{Axioms and Inference Rules}
	\begin{itemize}
		\item Certainty weakening:
		\[
		\text{If } \beta \leq \alpha, \quad (a, \alpha) \vdash (a, \beta)
		\]
		\begin{exampleblock}{Example}
		\((\text{It will rain tomorrow}, 0.9) \vdash (\text{It will rain tomorrow}, 0.7)\)  
		- If the certainty level decreases, the conclusion still holds.
		\end{exampleblock}
		\item Modus ponens:
		\[
		(\neg a \lor b, \alpha), \quad (a, \alpha) \vdash (b, \alpha)
		\]
		\begin{exampleblock}{Example}
			\((\neg \text{fever} \lor \text{flu}, 0.8), \quad (\text{fever}, 0.8) \vdash (\text{flu}, 0.8)\)  
			- If fever is observed, flu is inferred with certainty 0.8.
		\end{exampleblock}
		
	\end{itemize}
\end{frame}

\begin{frame}{Axioms and Inference Rules}
	\begin{itemize}
		\item Resolution:
		\[
		(\neg a \lor b, \alpha), \quad (a \lor c, \alpha) \vdash (b \lor c, \alpha)
		\]
		\begin{exampleblock}{Example}
			\((\neg \text{fever} \lor \text{flu}, 0.8), \quad (\text{fever} \lor \text{headache}, 0.8) \vdash (\text{flu} \lor \text{headache}, 0.8)\)
		\end{exampleblock}
		\item Weakest link resolution:
		\[
		(\neg a \lor b, \alpha), \quad (a \lor c, \beta) \vdash (b \lor c, \min(\alpha, \beta))
		\]
		\begin{exampleblock}{Example}
			\((\neg \text{fever} \lor \text{flu}, 0.9), \quad (\text{fever} \lor \text{headache}, 0.7) \vdash (\text{flu} \lor \text{headache}, 0.7)\)
			- The weakest certainty (0.7) determines the conclusion.
		\end{exampleblock}
	\end{itemize}
\end{frame}

\begin{frame}{Axioms and Inference Rules}
	\begin{itemize}
		\item Formula weakening:
		\[
		\text{If } a \vdash b \text{ in classical logic, then } (a, \alpha) \vdash (b, \alpha).
		\]
		\begin{exampleblock}{Example}
			\((\text{fever}, 0.9) \vdash (\text{flu}, 0.9)\)
			- Logical implication is preserved with the same certainty level.
		\end{exampleblock}
	\end{itemize}
\end{frame}

% The idea that in a reasoning chain, the certainty level of the conclusion is the smallest of the certainty levels of the formulas involved in the premises is at the basis of the syntactic approach proposed for plausible reasoning, and would date back to Theophrastus, an Aristotle?s follower.
% It turns out that any valid deduction in propositional logic is valid in possibilistic logic as well where the corresponding propositions are associated with any level ??(0,1]


% Slide 5: Handling Inconsistencies
\begin{frame}{Managing Inconsistent Information}
	\begin{itemize}
		\item Possibilistic logic can handle conflicting information by:
		\begin{itemize}
			\item Identifying an **inconsistency level** (\(inc-l\)).
			\item Filtering out rules with certainty levels below \(inc-l\).
		\end{itemize}
		\item **Example**:
		\[
		\{(a, 0.8), (\neg a, 0.7)\}
		\]
		\begin{itemize}
			\item Inconsistency level: \(inc-l = 0.7\).
			\item Filtered result: \((a, 0.8)\) dominates.
		\end{itemize}
	\end{itemize}
\end{frame}


% Slide 5: Logical Base and Weighted Clauses
\begin{frame}{Logical Base and Weighted Clauses}
	\begin{itemize}
		\item A possibilistic logic base is a collection of weighted formulas:
		\[
		\{(a_1, \alpha_1), (a_2, \alpha_2), \dots\}.
		\]
		\item Clauses can be rewritten in terms of their weights:
		\begin{itemize}
			\item Example:
			\[
			(a \land b, \alpha) \quad \Leftrightarrow \quad \{(a, \alpha), (b, \alpha)\}.
			\]
		\end{itemize}
		\item Logical tautologies (\(t\)) hold with any certainty level:
		\[
		(t, \alpha), \quad \forall \alpha \in [0, 1].
		\]
	\end{itemize}
\end{frame}

% Slide 6: Example Summary
\section{Example: Diagnoses}
\begin{frame}{Example: Medical Diagnosis}
	\begin{itemize}
		\item Rule base:
		\begin{align*}
			&(\neg \text{fever} \lor \text{flu}, 0.9), \\
			&(\neg \text{cough} \lor \text{cold}, 0.8), \\
			&(\text{fever} \lor \text{headache}, 0.7).
		\end{align*}
		\item Facts:
		\[
		(\text{fever}, 1), \quad (\text{cough}, 1).
		\]
		\item Inference:
		\begin{itemize}
			\item From Modus Ponens:
			\[
			(\text{flu}, 0.9), \quad (\text{cold}, 0.8).
			\]
			\item From Resolution:
			\[
			(\text{flu} \lor \text{headache}, 0.7).
			\]
		\end{itemize}
	\end{itemize}
\end{frame}

% Slide 7: Conclusion
\section{Conclusion}
\begin{frame}{Conclusion}
	\begin{itemize}
		\item Possibilistic logic provides a robust framework for reasoning under uncertainty.
		\item Key features:
		\begin{itemize}
			\item Certainty levels for qualitative inference.
			\item Compatibility with classical logic.
			\item Handling of conflicting information.
		\end{itemize}
		\item Applications:
		\begin{itemize}
			\item Medical diagnosis, decision-making, and knowledge representation.
		\end{itemize}
	\end{itemize}
\end{frame}

% Slide 2: Syntax of Possibilistic Logic
\section{Syntax and Semantics}
\begin{frame}{Basic Syntax}
	\begin{itemize}
		\item A possibilistic formula is a pair:
		\[
		(a, \alpha), \quad \text{where } \alpha \in [0, 1].
		\]
		\item Interpretation:
		\begin{itemize}
			\item \(a\): A propositional formula.
			\item \(\alpha\): Certainty level (degree of belief in \(a\)).
		\end{itemize}
		\item **Examples**:
		\begin{itemize}
			\item \((b \to f, 0.9)\): Birds fly with certainty 0.9.
			\item \((p \to \neg f, 1)\): Penguins do not fly, absolutely certain.
		\end{itemize}
	\end{itemize}
	\pause
	\begin{block}{Key Idea}
		Certainty levels enrich classical logic, enabling reasoning about typical cases and exceptions.
	\end{block}
\end{frame}

% Slide 3: Inference Rules
\begin{frame}{Key Inference Rules}
	\begin{itemize}
		\item **Certainty Weakening**:
		\[
		\underline{\text{If } \beta \leq \alpha, \quad (a, \alpha) \vdash (a, \beta)}.
		\]
		\item **Modus Ponens**:
		\[
		\underline{(\neg a \lor b, \alpha), \quad (a, \alpha) \vdash (b, \alpha)}.
		\]
		\item **Resolution**:
		\[
		\underline{(\neg a \lor b, \alpha), \quad (a \lor c, \alpha) \vdash (b \lor c, \alpha)}.
		\]
		\item **Weakest Link Resolution**:
		\[
		\underline{(\neg a \lor b, \alpha), \quad (a \lor c, \beta) \vdash (b \lor c, \min(\alpha, \beta))}.
		\]
	\end{itemize}
	\pause
	\begin{block}{Penguin Example}
		Consider:
		\begin{align*}
			(\neg b \lor f, 0.9), \quad (\neg p \lor \neg f, 1), \quad (\neg p \lor b, 1).
		\end{align*}
		If:
		\[
		(b, 1) \quad \text{(Tweety is a bird)}.
		\]
		Then:
		\[
		(\neg p \lor f, 0.9) \quad \text{(Tweety flies by default)}.
		\]
	\end{block}
\end{frame}

% Slide 4: Possibility Measures
\section{Possibility Measures}
\begin{frame}{Default Rules as Constraints}
	\begin{itemize}
		\item Default rule:
		\[
		\underline{\Pi(a \land b) > \Pi(a \land \neg b)}.
		\]
		- \(a \to b\): "If \(a\), then \(b\) is more plausible than \(\neg b\)."
		\item In terms of necessity:
		\[
		\underline{N(b | a) = 1 - \Pi(\neg b | a) > 0}.
		\]
		\item Comparative relation:
		\[
		\underline{a \land b >_\Pi a \land \neg b}.
		\]
	\end{itemize}
	\pause
	\begin{block}{Penguin Example Constraints}
		\begin{align*}
			b \land f >_\Pi b \land \neg f \quad &\text{(birds fly)}. \\
			p \land \neg f >_\Pi p \land f \quad &\text{(penguins don't fly)}. \\
			p \land b >_\Pi p \land \neg b \quad &\text{(penguins are birds)}.
		\end{align*}
	\end{block}
\end{frame}

% Slide 5: Computing Necessity Levels
\section{Example: Computing Necessity}
\begin{frame}{Computing Necessity for Penguin Example}
	\begin{itemize}
		\item **Possibility Distribution**:
		\[
		\pi(\omega) \text{ values assigned to worlds } \omega.
		\]
		\item Worlds:
		\[
		\Omega = \{ \omega_0, \omega_1, \omega_2, \omega_3, \omega_4, \omega_5, \omega_6, \omega_7 \}.
		\]
		\pause
		\item **Constraints**:
		\begin{align*}
			C_1 &: \max(\pi(\omega_6), \pi(\omega_7)) > \max(\pi(\omega_4), \pi(\omega_5)). \\
			C_2 &: \max(\pi(\omega_5), \pi(\omega_1)) > \max(\pi(\omega_3), \pi(\omega_7)). \\
			C_3 &: \max(\pi(\omega_5), \pi(\omega_7)) > \max(\pi(\omega_1), \pi(\omega_3)).
		\end{align*}
	\end{itemize}
\end{frame}

% Slide 6: Inference with Facts
\begin{frame}{Inference with New Information}
	\begin{itemize}
		\item Factual knowledge:
		\[
		(b, 1) \quad \text{(Tweety is a bird)}.
		\]
		\item Possibilistic inference:
		\[
		K \cup \{(b, 1)\} \vdash (f, 1/3).
		\]
		- Default rule: Tweety flies with certainty \(1/3\).
		\item Adding \(p = 1\) (Tweety is a penguin):
		\[
		K \cup F \vdash (\neg f, 2/3).
		\]
		- New inference: Tweety does not fly.
	\end{itemize}
\end{frame}

% Slide 7: Conclusion
\section{Conclusion}
\begin{frame}{Conclusion}
	\begin{itemize}
		\item Possibilistic logic effectively models:
		\begin{itemize}
			\item Default reasoning with exceptions.
			\item Reasoning under uncertainty.
		\end{itemize}
		\item **Penguin Example**:
		\begin{itemize}
			\item Illustrated how rules like "birds fly" and "penguins don't fly" can be reconciled.
			\item Demonstrated inference and conflict resolution.
		\end{itemize}
		\item **Applications**:
		\begin{itemize}
			\item Knowledge bases.
			\item Default rule modeling.
			\item Nonmonotonic reasoning.
		\end{itemize}
	\end{itemize}
\end{frame}
% Slide 6: Applications
\section{Applications}
\begin{frame}{Applications of Possibilistic Logic}
	\begin{itemize}
		\item **Medical Diagnosis**:
		\begin{itemize}
			\item Rules with certainty for differential diagnoses.
		\end{itemize}
		\item **Decision-Making**:
		\begin{itemize}
			\item Selecting actions under incomplete knowledge.
		\end{itemize}
		\item **Knowledge Bases**:
		\begin{itemize}
			\item Handling updates with belief revision.
		\end{itemize}
	\end{itemize}
	\begin{block}{Example Wrap-Up}
		Throughout the presentation, we used:
		\begin{itemize}
			\item Syntax: Representation of rules (\((a, \alpha)\)).
			\item Measures: Calculating \(\Pi(A)\) and \(N(A)\).
			\item Inference: Combining uncertain rules.
			\item Inconsistencies: Managing conflicts in rules.
		\end{itemize}
	\end{block}
\end{frame}

\begin{frame}{Applications of Possibility Theory}
	\begin{itemize}
		\item **Default Reasoning**: Handling incomplete knowledge in expert systems.
		\item **Medical Diagnosis**: Prioritizing plausible diagnoses based on symptoms.
		\item **Decision-Making**: Evaluating options under uncertainty.
		\item **Fuzzy Control Systems**: Representing and processing vague information.
	\end{itemize}
\end{frame}


\section{Introduction}
%*
%Handling Uncertainty in AI:
%Uncertainty is a common feature of real-world information and knowledge.
%Early AI systems, such as the MYCIN expert system (1970s), tackled uncertainty using certainty factors to represent degrees of belief and disbelief.
%2.	The Role of Possibility Theory:
%?	Possibility theory was introduced to represent imprecise and uncertain information. It complements probability theory by handling incomplete knowledge and degrees of plausibility.
%?	Developed by Zadeh (1978) and extended by others like Dubois and Prade, possibility theory provides a framework for qualitative reasoning about uncertainty.
%3.	Relation to Probability:
%?	Unlike probability, which requires exact numeric values, possibility theory uses dual measures:
%?	Possibility measure ($ \Pi $ ): Upper bounds of unknown probabilities.
%?	Necessity measure ( N ): Lower bounds, representing degrees of certainty.
%?	These measures allow for states of partial belief or ignorance, which are not well-handled by probabilistic approaches.
%4.	Possibilistic Logic:
%?	Possibilistic logic is the logical counterpart of possibility theory. It associates classical logic formulas with certainty levels, enabling qualitative reasoning.
%?	It can manage inconsistencies by stratifying formulas based on their certainty levels, prioritizing the most reliable information.

\begin{frame}{Introduction}

\end{frame}

\section{Syntax}

Possibility distribution: \\
$U$ is a set of affairs.\\
A possibility distribution is a mapping $ \pi $ from $ U $ to a totally ordered scale $ S $, with top denoted by 1 and bottom by 0. In the finite case $ S= \left\lbrace 1 = \lambda_{1} >...\lambda_{n}  >\lambda_{n+1}  = 0 \right\rbrace  $. \\
The function $ \pi $ represents the state of knowledge of an agent (about
the actual state of affairs), also called an epistemic state.\\
$ \pi \left( u \right) = 0 $  means that state $ u $ is rejected as impossible;\\
$ \pi \left( u \right) = 1 $ means that state $ u $ is totally possible (= plausible).\\
Specificity:
A possibility distribution $ \pi $ is said to be at least as specific
as another $ \pi ' $ if and only if for each state of affairs $ u $: $ \pi \left( u \right) \leq \pi' \left( u \right) $. Then, $ \pi $ is at least as restrictive and informative as  $ \pi ' $ , since it
rules out at least as many states with at least as much strength.\\
Complete knowledge: for some $ u_{0} $, $ \pi \left( u_{0} \right) = 1 $ and $ \pi \left( u \right) = 0 $, $ \forall u  \neq u_{0}$ (only $ u_{0} $, is possible). \\
Complete ignorance: $ \pi \left( u \right) = 1 $,$ \forall u  \in U $ (all states are possible).

Principle of minimal specificity: It states that any hypothesis not known to be impossible cannot be ruled out. It is a minimal commitment, cautious, information principle.

Possibility and necessity functions:
$\Pi(A) = \sup_{u \in A} \pi(u); \quad N(A) = \inf_{u \notin A} (1 - \pi(u))$



\section{Section}

\begin{frame}{Blocks}
    \begin{block}{Block}
        This is a regular block.
        \begin{itemize}
            \item This is an item in a block.
        \end{itemize}
    \end{block}
    
    \begin{alertblock}{Block}
        This is an alert block.
        \begin{itemize}
            \item This is an item in an alert block.
        \end{itemize}
    \end{alertblock}
    
    \begin{exampleblock}{Block}
        This is an example block.
        \begin{itemize}
            \item This is an item in an example block.
        \end{itemize}
    \end{exampleblock}
\end{frame}


\subsection{Subsection}

\begin{frame}{Highlighting Text}
    We can \highlight{highlight} text. There are also a \darkhighlight{darker} option, and a \yellowhighlight{yellow} option.\\
    The same options are also available in math mode: $\mathhighlight{a} + \darkmathhighlight{b} = \yellowmathhighlight{c}$
\end{frame}


\appendix

\section{Appendix}

\begin{frame}{Appendix}
    This is an appendix.
    \begin{itemize}
        \item Page numbers start from the beginning.
    \end{itemize}
\end{frame}

\end{document}
